\section{Protocole d'Évaluation (Benchmarking)}

Le développement d'une architecture agentique complexe ne constitue que la moitié du travail d'ingénierie. L'autre moitié, tout aussi cruciale, réside dans la capacité à mesurer objectivement la performance du système. L'évaluation automatique d'un assistant IA est un domaine notoirement difficile, car contrairement à des tâches de classification, il n'existe pas toujours une unique "bonne réponse" syntaxique pour un problème donné.

\subsection{Le défi de la métrique : L'illusion de la similarité}

Dans les systèmes de NLP classiques, l'étalon-or est souvent la \textit{Cosine Similarity} (similarité cosinus). Cette métrique géométrique calcule l'angle entre le vecteur de la réponse générée et celui de la réponse de référence (Vérité Terrain). Si l'angle est faible, les textes sont considérés comme sémantiquement proches.

Cependant, dans le contexte rigoureux d'un DSL comme Envision, cette approche s'est révélée imprécise, voire contre-productive : un modèle peut générer une réponse syntaxiquement parfaite, utilisant le bon vocabulaire technique, ce qui lui confère un score de similarité très élevé. Pourtant, une simple négation dans une phrase ou une référence à la mauvaise table (\texttt{Orders} au lieu de \texttt{OrderLines}) rend la réponse objectivement fausse. La métrique vectorielle ne capture pas cette "rupture" logique.

\subsection{La solution polyvalente : "LLM-as-a-Judge"}

Pour pallier ces biais structurels, nous avons implémenté une méthode d'évaluation moderne basée sur un \textbf{Juge Artificiel}, inspirée des travaux récents sur l'évaluation des chatbots \cite{Zheng2023}. 

Le principe est de déléguer l'évaluation à un LLM tiers (dans notre cas, un modèle comme \texttt{mistral} ou \texttt{deepseek-v3}) dont l'unique tâche est de comparer la réponse de l'agent avec la vérité terrain. Il analyse la cohérence logique et factuelle de la réponse. Contrairement à la solution "cosine similarity", on garantit qu'aucune réponse fausse mais sémantiquement similaire (en terme d'\textit{embedding}) ne soit valorisée et on récompense les réponses justes et appropriées. 

\subsection{L'utilisation du dual cross-encoder}

Comme expliqué précédement dans ce rapport, un \textit{cross-encoder} est un modèle entraîné pour comparer directement deux documents et donner un score de similarité. 
\\

Concrètement, le cross-encoder utilise une architecture de transformer. Il concatène les deux documents grâce à un séparateur spécial ce qui permet aux tokens des deux passages de s'influencer entre eux. 
\\

L'approche duale combine deux perspectives. Tout d'abord, la cohérence logique (NLI - Natural Language Interference, modèle : DeBERTa v3) : on évalue si la réponse générée est logiquement cohérente avec la réponse attendue ou si des contradictions factuelles existent. En quelque sorte, il s'agit de garantir la justesse factuelle. Ensuite, la pertinence contextuelle (modèle : Jina): on cherche à mesurer le degré d'alignement entre la question posée et la réponse générée. Il s'agit d'observer la pertinence et l'utilité de la réponse. 
\\

Enfin le modèle vérifie d'abord s'il n'y a pas de contradiction majeure, auquel cas, le score sera zéro. Sinon, il combine les deux critères. Ce benchmark a l'avantage de fournir un scoring continu contrairement au \textit{LLM-as-a-Judge} mais est beaucoup plus compliqué à faire fonctionner en pratique du fait de son grand nombre d'hyper-paramètres à ajuster (Quel est le seuil d'une contradiction majeure ? Comment combiner les 2 critères ? etc).

\subsection{La solution hybride}

Les deux solutions présentées ci-dessus permettent de mieux noter les réponses. En effet, la réponse proposée, même si elle est proche en termes géométriques, n'obtiendra pas une bonne note si elle est incorrecte. Cependant le problème qui se pose est celui de l'utilisation d'un modèle (LLM ou Cross-Encoder), souvent coûteux, afin de vérifier que la réponse est juste.

Or, lorsqu'il s'agit de retourner un nombre précis ou une liste de fichiers exhaustive pour une question structurelle ou syntaxique, la notation attendue est binaire: "\textit{les éléments attendus sont-ils présents ou non ?}". Nous pouvons donc simplement faire une recherche par regex dans la réponse produite par notre architecture, pour vérifier que la donnée précise y est contenue. Cette approche a de plus l'avantage d'être infaillible sur les questions structurelles contrairement aux approches reposant sur les transformers qui sont plus heuristiques.
\\

Ainsi, nous pouvons facilement et efficacement corriger les questions structurelles, et faire le choix d'une correction par la méthode \textit{LLM-as-a-judge} ou \textit{Dual Cross-Encoder} pour les questions sémantiques.

\subsection{Constitution du Benchmark}
Pour alimenter le benchmark, nous avons construit manuellement un jeu de données de test composé de paires questions-réponses issues de l'historique réel de développement de Lokad. Ce dataset est segmenté en deux catégories distinctes pour tester les différentes capacités de notre architecture hybride :

\begin{enumerate}
    \item \textbf{Questions structurelles/syntaxiques (Test du GREP) :} 
    \\ \textit{Exemple :} "Quels scripts lisent le fichier \texttt{/Clean/Items.ion} ?" 
    \\ \textit{Attendu :} Une liste exhaustive et exacte des noms de fichiers. Ici, l'évaluation est binaire : la liste est complète ou elle ne l'est pas. C'est un test de la capacité de l'agent à naviguer dans la structure du projet.
    
    \item \textbf{Questions sémantiques (Test du RAG) :} 
    \\ \textit{Exemple :} "Où est définie la logique d'actualisation du stock ?" 
    \\ \textit{Attendu :} Une réponse pointant vers les bons modules et expliquant brièvement la logique. Ici, le Juge évalue la pertinence des documents retrouvés et la clarté de la synthèse.
\end{enumerate} 

\subsection{Métriques de performance système}

Au-delà de la qualité intrinsèque des réponses (précision), la viabilité industrielle de l'outil dépend de sa performance opérationnelle. Nous avons donc instrumenté notre pipeline pour suivre deux indicateurs critiques :

\begin{itemize}
    \item \textbf{La Latence d'Inférence :} 
    Nous nous intéressons au temps écoulé entre la requête utilisateur et la réponse finale. Actuellement, une boucle agentique complète (réflexion, appel d'outil, génération) prend entre 5 et 20 secondes. C'est acceptable pour un assistant de développement, mais une piste d'amélioration consiste à réduire ce délai en optimisant les appels réseaux.
    
    \item \textbf{Le Temps de Construction de la base de données :} 
    Une métrique souvent négligée mais vitale. Le code source d'Envision évolue tous les jours. Si notre système met 30 minutes à "lire" et indexer la base de code, il sera toujours obsolète.
    \\ Grâce à notre parser optimisé, nous parvenons actuellement à parser, chunker et indexer l'ensemble du corpus (plusieurs centaines de scripts) en moins de \textbf{1 minute}.
\end{itemize}

\subsection{Evaluation du benchmark}

Pour vérifier la qualité du benchmark, nous avons implémenté un benchmark “pour le benchmark” : nous avons généré un certain nombre de réponses aux questions du benchmark, complété par des versions modifiées de ces réponses pour avoir des exemples de réponses correctes et incorrectes pour chaque question. Pour chacune de ces réponses, nous avons indiqué si elle était correcte ou non, puis nous avons lancé le benchmark et comparé les notes du benchmark avec nos indications. \textit{LLM-as-a-judge} obtient un score entre 19/20 et 20/20 sur ce “benchmark pour benchmark”, confirmant que notre benchmark est plutôt fiable.